\documentclass[10pt,a4paper]{amsart}
\usepackage[margin=19mm]{geometry}
\usepackage{amsmath,amssymb,mathtools}
\usepackage[scr=rsfs,cal=euler]{mathalfa}
\usepackage{xfrac}
\usepackage[unicode,hidelinks,bookmarks=false]{hyperref}
\newcommand{\ie}{i.e.\ }
\newcommand{\cf}{cf.\ }
\newcommand{\resp}{respectively\ }
\newcommand{\Subsection}{\subsection}
\providecommand{\nobreakdash}{\nobreak\hbox{-}}
\setlength{\emergencystretch}{2em}
\multlinegap0pt
\usepackage{bm}
\usepackage{bbm}
\usepackage{dsfont}
\usepackage{xspace}
\usepackage{cancel}
\usepackage{mathtools}
\usepackage{cleveref}
\usepackage{amsthm}
\makeatletter
\@ifundefined{theorem}{\newtheorem{theorem}{Theorem}}{}
\@ifundefined{lemma}{\newtheorem{lemma}[theorem]{Lemma}}{}
\@ifundefined{proposition}{\newtheorem{proposition}[theorem]{Proposition}}{}
\makeatother
\title[On positive automorphisms of algebras of operators on atomic]{On positive automorphisms of algebras of operators on atomic Archimedean vector lattices}
\author{Gregor Cigler, Marko Kandi\'c}
\date{Source: arXiv:2601.21811v1 (CC BY 4.0)}
\begin{document}
\maketitle

\noindent\emph{Source and licence.} On positive automorphisms of algebras of operators on atomic Archimedean vector lattices. Version arXiv:2601.21811v1 (CC BY 4.0), distributed under the Creative Commons Attribution 4.0 International licence, as recorded in the arXiv licence metadata for this exact version and shown on the abstract page https://arxiv.org/abs/2601.21811v1. Full rights notice: licenses/arxiv-2601.21811-CC-BY-4.0.txt.\par\smallskip

\noindent\textit{Editorial scope.} Counted unit: the Theorem whose source label is \texttt{main theorem} in the article (source0.tex lines 533--546, source label \texttt{main theorem}), delivered together with the article's own complete original proof of it (source0.tex lines 548--589, one \texttt{proof} environment of 42 lines). No mathematics is written, repaired or re-derived: statement and proof are the article's own text line for line. Required context retained uncounted: zapis z atomi (lines 196--198); nicelni operator atomi (lines 276--284); A\_0 mreza (lines 355--362); lema o delti (lines 437--446); a0 pride iz a0 (lines 514--516). Every definition, display and label of those lines is retained unchanged. Omitted from this package and not redistributed: 1--195, 199--275, 285--354, 363--436, 447--513, 517--532, 547--547, 590--900. Nothing inside a retained excerpt is omitted or altered: every retained body is delivered exactly as the article's lines are written, so no omission receipt of any kind accompanies this package. Citations: the retained excerpts carry no citation command at all, so no bibliography entry of the article is transcribed and no reference is redistributed. Reference adaptation: none was needed and none was made; every reference inside the delivered unit resolves inside it, so no unresolved reference or citation is produced. Printed slips kept literal: no printed word, symbol, hypothesis or number of the counted statement or of its proof was altered. Layout and class: the article's own class is \texttt{amsart} with packages including hyperref, cleveref, amssymb, eqnarray, mathrsfs, xcolor, amsmath; the wrapper is the lane's pinned \texttt{amsart} editorial wrapper at 10pt on A4, which restates the theorem-like environments the retained text needs and adds the article's own macro definitions that the retained text needs (none), with extra wrapper packages (bm, bbm, dsfont, xspace, cancel, mathtools, cleveref). The delivered numbering is the unit's own. Assets: no retained span contains a figure environment, a graphics-inclusion command, a PDF-page inclusion, a table or a TikZ picture; no image file is copied into this package, the package declares no asset dependency and no image file is redistributed with it. Licence: version arXiv:2601.21811v1 of this article is distributed under the Creative Commons Attribution 4.0 International licence, as recorded in the arXiv licence metadata for this exact version and shown on the abstract page https://arxiv.org/abs/2601.21811v1. A full rights notice is delivered at licenses/arxiv-2601.21811-CC-BY-4.0.txt. Characters: every retained excerpt is pure ASCII, so the delivered engine, XeLaTeX with its default Latin Modern text font, reports no missing character.
\begin{equation}\label{zapis z atomi}
x=\sup_{a\in \mathcal A}\varphi_a(x)a=\sup_{F\in \mathcal F}\sup_{a\in F}\varphi_a(x)a=\sup_{F\in \mathcal F}\sum_{a\in F}\varphi_a(x)a. 
\end{equation}

\begin{lemma}\label{nicelni operator atomi}
Let $T\colon X\to Y$ be a linear operator between Archimedean vector lattices and let $\mathcal A$ and $\mathcal B$ be maximal families of pairwise disjoint atoms in $X$ and $Y$, respectively.
\begin{enumerate}
\item If $X$ is atomic and $T$ is order continuous, then $T\geq 0$ if and only if $Ta\geq 0$ for every $a\in \mathcal A$.  
\item If $Y$ is atomic, then $T\geq 0$ if and only if $\varphi_b\circ T\geq 0$ for every atom $b\in \mathcal B$.
\item If $X$ and $Y$ are atomic and $T$ is order continuous, then $T\geq 0$ if and only if $\varphi_b(Ta)\geq 0$ for all $a\in \mathcal A$
 and $b\in \mathcal B$.
\end{enumerate}
\end{lemma}

\begin{proposition}\label{A_0 mreza}
    For an Archimedean vector lattice $X$ the following statements hold.
    \begin{enumerate}
        \item The algebra $\mathscr A_0$ is a vector lattice. For $T=\sum_{a,b\in F}\lambda_{ab}(a\otimes \varphi_b)\in \mathscr A_0$ the modulus is given by 
        $$|T|=\sum_{a,b\in F}|\lambda_{ab}|(a\otimes \varphi_b).$$
        \item If $X$ is Dedekind complete, then $\mathscr A_0$ is a vector sublattice of $\mathscr L_n(X)$.
    \end{enumerate}
\end{proposition}

\begin{lemma}\label{lema o delti}
If $\Phi(\mathscr A_0)\subseteq \mathscr L_n(X)$, then the following assertions hold.
\begin{enumerate}
\item For all $a,b\in \mathcal A$ the operator $F_{ab}$ is a rank-one operator.
\item For all $a,b\in \mathcal A$ there exist positive scalars $\gamma_a$, a positive nonzero vector $u_a\in X$ and a positive nonzero order continuous functional $\psi_b$ such that
    $$F_{ab}=\tfrac{\gamma_a}{\gamma_b} (u_a\otimes \psi_b)$$
    with $\psi_a(u_a)=1$ and $\psi_a(u_b)=0$ for $b\neq a$.
\item For distinct atoms $a,b\in \mathcal A$ and an atom $e$ we have $\psi_b(e)=0$ whenever $u_a\wedge e >0$.
\end{enumerate}
\end{lemma}

\begin{proposition}\label{a0 pride iz a0}
Let $\mathscr A\subseteq \mathscr L_n(X)$ be any algebra of operators which contains $\mathscr A_0$. If $\Phi$ is a positive automorphism of $\mathscr A$, then for arbitrary $T\in \mathscr A$ we have $T\in \mathscr A_0$ whenever $\Phi(T)\in \mathscr A_0$.
\end{proposition}

\begin{theorem}\label{main theorem}
Let $\mathscr A\subseteq \mathscr L_n(X)$ be any algebra of operators that contains $\mathscr A_0$. The following assertions hold for a positive algebra automorphism $\Phi\colon \mathscr A\to \mathscr A$.
\begin{enumerate}
  \item There exist a bijection $\pi\colon\mathcal A\to\mathcal A$ and a set of positive scalars $\{\delta_a:\; a\in\mathcal A\}$ such that
for all $a,b\in\mathcal A$ we have
\begin{align}
\Phi(a\otimes\varphi_b)&=\tfrac{\delta_a}{\delta_b} \left(\pi(a)\otimes\varphi_{\pi(b)}\right) \label{formula1}\\
\Phi^{-1}(a\otimes\varphi_b)&=\tfrac{\delta_{\pi^{-1}(b)}}{\delta_{\pi^{-1}(a)}} \left(\pi^{-1}(a)\otimes\varphi_{\pi^{-1}(b)}\right). \label{formula2}
\end{align}
In particular,
$\Phi(\mathscr A_0)=\mathscr A_0$ and $\Phi(\mathscr A_0^+)=\mathscr A_0^+$. 
  \item The inverse $\Phi^{-1}$ is positive.
\end{enumerate}
\end{theorem}

\begin{proof}
(i) First we prove that for arbitrary $a\in\mathcal A$ the set ${\mathcal U}_a$ is a singleton set.
Assume on the contrary that for some $a$ the set $\mathcal U_a$ contains two different atoms $e$ and $f$, and pick $b, c\in \mathcal A$. Then 
$$F_{bc}e=\frac{\gamma_b}{\gamma_c}(u_b\otimes\psi_c)(e)=\frac{\gamma_b}{\gamma_c}\psi_c(e)u_b.$$
If $c\ne a$, by \Cref{lema o delti}(iii) we conclude $\psi_c(e)=0$ which implies
$F_{bc}e=0$, and similarly, $F_{bc}(f)=0$.

By \Cref{a0 pride iz a0} there exists a finite set $F\subseteq \mathcal A$ which without loss of generality contains $a$ such that $\Phi(T)=e\otimes \varphi_e +f\otimes\varphi_f$  where $T=\sum_{b,c\in F}t_{bc}(b\otimes\varphi_c)\in \mathscr A_0$. Then 
$$\Phi(T)=\sum_{b,c\in F}t_{bc}F_{bc}.$$
Since for $c\ne a$ we have $F_{bc}e=0$ and $F_{ba}e=\frac{\gamma_b}{\gamma_a}\psi_a(e)u_b$ it follows that
$$e=\Phi(T)e=\sum_{b,c\in F}t_{bc}F_{bc}e=\sum_{b\in F}t_{ba}F_{ba}e= \tfrac{\gamma_b}{\gamma_a}\psi_a(e)\sum_{b\in\mathcal F}t_{ba}u_b.$$
Similarly we get
$$f=\tfrac{\gamma_b}{\gamma_a}\psi_a(f)\sum_{b\in F}t_{ba}u_b.$$
which is a contradiction since $e$ and $f$ are linearly independent. This proves that for each $a\in \mathcal A$ the set $\mathcal U_a$ is a singleton set which yields an injective mapping $\pi\colon \mathcal A\to \mathcal A$ such that ${\mathcal U}_a=\{\pi(a)\}$. 

Since $\mathcal U_a$ consists of all atoms that are not disjoint with $u_a$ and $X$ is atomic, an application of \eqref{zapis z atomi} yields $u_a=\kappa_a\pi(a)$ for some positive scalar $\kappa_a$. Hence, 
$$\Phi(a\otimes \varphi_b)=\kappa_a\pi(a)\otimes \psi_b.$$
To prove that $\pi$ is surjective, choose an atom $c\in\mathcal A$. Since $\Phi$ is an automorphism of $\mathscr A$, there exists $T\in \mathscr A$ such that $\Phi(T)=c\otimes \varphi_c$. By
\Cref{a0 pride iz a0} we conclude $T\in \mathscr A_0$. Therefore, there exists a finite set $F\subseteq \mathcal A$ such that 
$T=\sum_{a,b\in F} \lambda_{ab}(a\otimes \varphi_b)$ for some scalars $\lambda_{ab}$.  Since
$$c=(c\otimes \varphi_c)c=\Phi(T)c=\sum_{a,b\in F}\kappa_a (\pi(a)\otimes \psi_b)c=\sum_{a,b\in F}\kappa_a \psi_b(c)\pi(a),$$ and $\pi(a)\in \mathcal A$ for each $a\in F$, we conclude that $c=\pi(a)$ for some $a\in \mathcal A$. 

For different atoms $a,b\in\mathcal A$ we have $\kappa_a\psi_b(\pi(a))=\psi_b(u_a)=0$. Since $\kappa_a\ne 0$ we get $\psi_b(\pi(a))=0$ for all $a\ne b$. Order continuity of $\psi_b$ implies that
$$\psi_b=\eta_b\varphi_{\pi(b)}$$
for some $\eta_b>0$, and so, finally
$\Phi(a\otimes\varphi_b)=\kappa_a\eta_b(\pi(a)\otimes\varphi_{\pi(b)}).$
Since $(\Phi(a\otimes\varphi_a))^2=\Phi(a\otimes\varphi_a)$ it follows that $\kappa_a\eta_a=1$, yielding $\eta_a=\frac{1}{\kappa_a}$.
By \Cref{lema o delti} we conclude
$$\Phi(a\otimes\varphi_b)=\tfrac{\gamma_a}{\gamma_b} (u_a\otimes \psi_b)=\tfrac{\gamma_a\kappa_a}{\gamma_b\kappa_b}(\pi(a)\otimes\varphi_{\pi(b)}).$$
By introducing $\delta_a=\gamma_a\kappa_a$ for each $a\in \mathcal A$ we obtain \eqref{formula1}. One can verify \eqref{formula2} by a direct calculation. 

To prove $\Phi(\mathscr A_0)=\mathscr A_0$ and $\Phi(\mathscr A_0^+)=\mathscr A_0^+$, note that \Cref{A_0 mreza}(i), and \eqref{formula1} and \eqref{formula2} imply
$\Phi(\mathscr A_0^+)\subseteq  \mathscr A_0^+$ and $\Phi^{-1}(\mathscr A_0^+)\subseteq \mathscr A_0^+$, respectively, yielding $\Phi(\mathscr A_0^+)=\mathscr A_0^+$. Since $\mathscr A_0=\mathscr A_0^+-\mathscr A_0^+$, we obtain $\Phi(\mathscr A_0)=\mathscr A_0$.

(ii) Pick a positive operator $T$ in $\mathscr A$ and atoms $a,b\in \mathcal A$. By \eqref{formula2}, there exist atoms $c,d\in \mathcal A$ such that
$\Phi^{-1}(c\otimes \varphi_c)=b\otimes \varphi_b$ and $\Phi^{-1}(d\otimes \varphi_d)=a\otimes \varphi_a$.
Since $(c\otimes \varphi_c)T(d\otimes \varphi_d)=\varphi_c(Td)(c\otimes \varphi_d) \in \mathscr A_0^+$ and by (i) we have $\Phi^{-1}(\mathscr A_0^+)=\mathscr A_0^+$, we conclude that
\begin{align*}
\Phi^{-1}((c\otimes \varphi_c)T(d\otimes \varphi_d))a&=(b\otimes \varphi_b)\Phi^{-1}(T)(a\otimes \varphi_a)a=\varphi_b(\Phi^{-1}(T)a)b
\end{align*}
is a positive vector. In particular, $\varphi_b(\Phi^{-1}(T)a)\geq 0$ for all $a,b\in \mathcal A$. To finish the proof we apply the fact that $\Phi^{-1}(T)$ is order continuous and \Cref{nicelni operator atomi}. 
\end{proof}

\end{document}
